\section{\silobench{}}
\label{sec:benchmark}

This section presents the architecture of \silobench{}, a configurable environment for evaluating multi-agent collaboration under information silos. 
Each configuration is defined by three orthogonal dimensions: agent scale $N$, communication protocol $P$, and language model $\mathcal{M}$. We describe the task space, evaluation metrics, and execution pipeline.

\subsection{Task Space}
\label{sec:task_space}

A central design goal of \silobench{} is to ground task difficulty in principled communication complexity theory, so that observed performance gaps can be attributed to coordination demands rather than ad hoc task choice. 
The theoretical foundation for analyzing distributed computation costs dates back to Yao's seminal work on communication complexity \citep{yao1979some}, which established the framework for quantifying the minimum bits required for distributed parties to compute a function. 
Building on this foundation, we categorize tasks by their optimal communication complexity:

\begin{equation}
\tau_k = (f_k, \mathcal{X}_k, y^*_k)
\label{eq:task_def}
\end{equation}
where $f_k$ specifies the computational function, $\mathcal{X}_k$ is the global input data, and $y^*_k$ is the ground-truth answer. 
Tasks are organized into three levels based on their optimal communication complexity (complete task specifications are provided in Appendix~\ref{app:tasks}).

\paragraph{Level I: Aggregation ($\mathcal{O}(N)$ communication).}
As illustrated in Figure~\ref{fig:task_levels} (left), these tasks exhibit embarrassingly parallel structure followed by reduction. 
Each agent processes its local shard independently, producing intermediate results aggregated through associative operations (e.g., \texttt{max}, \texttt{sum}, \texttt{xor}). 
The optimal topology is a star or tree structure where one agent collects all partial results. 
Representative tasks include global maximum (LC-414: ``Third Maximum Number''), distributed voting (LC-169: ``Majority Element''), and word frequency counting (LC-2085).

\paragraph{Level II: Mesh Network ($\mathcal{O}(N)$ communication).}
As shown in Figure~\ref{fig:task_levels} (center), these tasks exhibit spatial locality: agent $i$'s computation depends primarily on neighboring agents $i-1$ and $i+1$. 
Information propagates through a structured mesh via pairwise exchanges, with optimal topology being a linear chain requiring $N-1$ point-to-point exchanges. 
Representative tasks include prefix sum (LC-1480), moving average (LC-346: ``Moving Average from Data Stream''), and trapping rain water (LC-42).

\paragraph{Level III: Global Shuffle ($\mathcal{O}(N \log N)$ to $\mathcal{O}(N^2)$ communication).}
As depicted in Figure~\ref{fig:task_levels} (right), these tasks feature irregular, potentially all-to-all communication patterns where any agent's output may depend on information from any other agent. 
The range $\mathcal{O}(N \log N)$--$\mathcal{O}(N^2)$ spans from the classical lower bound for distributed reorganization to the full-consensus cost imposed by our evaluation criterion, where every agent must output the complete global answer. 
Representative tasks include distributed sorting (LC-912), graph connectivity (LC-323), and matrix multiplication (LC-311).


\subsection{Task Construction Pipeline.}
LeetCode problems serve solely as algorithmic inspiration and we \textit{do not} transform raw LeetCode data. 
For each task \emph{category} (e.g., ``Global Maximum''), we independently implement a Python generator that programmatically produces random inputs and exact ground-truth answers. 
A \emph{task instance} is one concrete input--output pair drawn from this generator under a fixed $(N, P, \mathcal{M})$ configuration, where $N$ is the agent scale, $P$ the communication protocol, and $\mathcal{M}$ the language model, and a fixed random seed, ensuring reproducibility while allowing unlimited fresh instances.

To illustrate: for Level-I Global Maximum (inspired by LC-414, the ``Third Maximum Number'' problem), given agent count $N$ and per-agent shard size $k$, the generator (i) samples $N \times k$ integers uniformly at random, (ii) partitions them into $N$ equal shards $\mathcal{X}_1,\ldots,\mathcal{X}_N$, (iii) computes $y^* = \max(\mathcal{X})$, and (iv) records the fixed seed for reproducibility. 
Each agent receives only its local shard $\mathcal{X}_i$ and must coordinate to determine $y^*$. 
This pipeline generalizes directly to all 30 tasks: the generator encodes the task-specific function $f_k$, scales global input size proportionally with $N$ to maintain constant per-agent workload, and produces exact ground-truth answers enabling fully objective evaluation.



\subsection{Evaluation Metrics}
\label{sec:metrics}

We define four complementary metrics to capture both \textit{what} agents achieve and \textit{how} they coordinate. 
Let $\hat{y}_i$ denote agent $i$'s submitted answer, and let $m_i$ denote the total number of messages successfully transmitted outward by agent $i$ during the entire collaboration.

\paragraph{Success Rate ($\mathcal{S}$).}
Measures the proportion of agents converging to the correct answer:
\begin{equation}
\mathcal{S} = \frac{1}{N} \sum_{i=1}^{N} \mathbbm{1}[\hat{y}_i = y^*]
\end{equation}
A task instance is \textit{successful} when $\mathcal{S} = 1$, indicating unanimous convergence.

\paragraph{Partial Correctness Score ($\mathcal{P}$).}
Binary success rate can understate partial progress. 
We introduce a continuous measure of answer quality tailored to each task category: for Level-I, $\mathcal{P}$ is the fraction of agents within tolerance of the ground truth; for Level-II, the fraction of correctly computed elements per local segment; for Level-III, the longest correctly ordered subsequence relative to total length. 
Letting $q_i \in [0,1]$ denote the per-agent quality score:
\begin{equation}
\mathcal{P} = \frac{1}{N} \sum_{i=1}^{N} q_i
\end{equation}
Together with $\mathcal{S}$, this score allows us to isolate where coordination breaks down: the gap $\mathcal{P} - \mathcal{S}$ quantifies performance lost specifically at the reasoning-integration stage rather than at the communication stage.

\paragraph{Token Consumption ($\mathcal{C}$).}
Quantifies computational cost per communication round:
\begin{equation}
\mathcal{C} = \frac{\sum_{i=1}^{N} \sum_{r=1}^{R} t_i^{\text{out}}[r]}{R_{max}}
\end{equation}
where $t_i^{\text{out}}[r]$ is the number of output tokens generated by agent $i$ in round $r$, and $R_{max}$ is the max number of rounds executed.

\paragraph{Communication Density ($\mathcal{D}$).}
Captures inter-agent interaction intensity. 
Here $N(N-1)$ is the directed-edge count when each ordered pair exchanges exactly one message; since agents may send multiple messages to the same recipient across rounds, $\mathcal{D} \in [0, +\infty)$:
\begin{equation}
\mathcal{D} = \frac{\sum_{i=1}^{N} m_i}{N(N-1)}
\end{equation}
Values near 0 suggest sparse, targeted exchanges; $\mathcal{D} = 1$ indicates one message per directed pair on average; values exceeding 1 reflect iterative multi-round exchanges. For the SFS protocol (see Appendix~\ref{app:protocols}), $m_i$ counts the number of times other agents successfully read files written by agent $i$, preserving the same ``information actually transferred'' semantics as direct message-passing.

Together, $\mathcal{S}$ and $\mathcal{P}$ measure \textit{what} agents achieve, $\mathcal{C}$ measures \textit{at what cost}, and $\mathcal{D}$ reveals \textit{how} they coordinate.

\subsection{Execution Pipeline}
\label{sec:pipeline}

Given task $\tau = (f, \mathcal{X}, y^*)$ and configuration $(N, P, \mathcal{M})$, the evaluation proceeds through four phases.

\paragraph{Phase 1: Data Partition.}
$\textsc{Partition}(\mathcal{X}, N) \rightarrow \{\mathcal{X}_1, \ldots, \mathcal{X}_N\}$, where $|\mathcal{X}_i| \approx |\mathcal{X}|/N$ ensures equal partition and no agent holds privileged information.

\paragraph{Phase 2: Agent Initialization.}
Each agent $i$ is initialized with model $\mathcal{M}$ and receives $\textsc{Init}(i) \leftarrow (\text{desc}(f, \mathcal{X}_i), P)$, specifying the core task logic, local data, and protocol constraint. 
The prompt provides task-structural guidance while preserving strategic autonomy (see Appendix~\ref{app:prompt_design}).

\paragraph{Phase 3: Collaborative Execution.}
Agents engage in iterative communication for up to $R_{\max}$ rounds. 
All $N$ agents are activated in parallel within each round: they receive incoming messages from the previous round, independently decide on actions, and execute them simultaneously. 
Messages or files written in round $r$ become visible at the start of round $r+1$. 
Execution terminates when all agents submit answers or the round limit is reached.

\paragraph{Phase 4: Metric Computation.}
The four metrics $(\mathcal{S}, \mathcal{P}, \mathcal{C}, \mathcal{D})$ are computed from submitted answers $\{\hat{y}_i\}_{i=1}^{N}$ and recorded communication logs.